\clearpage
\begin{explanationbox}
\textbf{\textcolor{CExplanation}{一、为什么会想到这条公式？从“转半圈”倒推}}
\begin{center}\begin{tikzpicture}[font=\small,>=stealth,line join=round]
\draw[gray!40,rounded corners=3pt] (0,0) rectangle (14.2,8.3);\node at (7.1,7.700000000000001) {先看一对点：中心就是连线的中点};
\begin{scope}[shift={(1.3,1.5)},x=1.4cm,y=1.1cm]\draw[->,gray] (0,0)--(4.7,0) node[right]{$x$};\draw[->,gray] (0,0)--(0,4.8) node[above]{$y$};\draw[blue,thick] (1,2)--(3,4);\fill[orange!85!black] (1,2) circle(2.4pt);\fill[orange!85!black] (3,4) circle(2.4pt);\fill[violet] (2,3) circle(2.5pt);\node[below right=7pt] at (1,2){$Q(1,2)$};\node[right=7pt] at (3,4){$P(3,4)$};\node[below right=8pt,violet] at (2,3){$C(2,3)$};\end{scope}\node[anchor=west,align=left] at (8.4,4.5){横坐标平均为2\\[14pt]纵坐标平均为3\\[14pt]$\overrightarrow{CP}=-\overrightarrow{CQ}$};
\node[text width=12.6cm,align=center] at (7.1,.65) {不指定整条函数曲线；一对点的对称，不能证明整张图像对称。};
\end{tikzpicture}\end{center}



先不要记公式。把一张画有曲线的透明纸，用大头针固定在 $C(2,3)$，转动 $180^\circ$。
如果曲线仍与原曲线重合，就是关于 $C$ 中心对称。
例如原来曲线上有 $P(3,4)$，转半圈后应落到 $Q(1,2)$。
\[
\underbrace{3+1}_{\text{横坐标和}}=4=2\times2,\qquad
\underbrace{4+2}_{\text{纵坐标和}}=6=2\times3.
\]
由此发现：\textbf{对称的两点，横坐标的中间数是 $a$，纵坐标的中间数是 $b$。}
这就是初中学过的\textbf{线段中点}，并不需要先背新公式。



\medskip

\end{explanationbox}
\clearpage
\begin{explanationbox}
\textbf{\textcolor{CExplanation}{二、怎样从中点“长出” $f(a+t)+f(a-t)=2b$？}}
\begin{center}\begin{tikzpicture}[font=\small,>=stealth,line join=round]
\draw[gray!40,rounded corners=3pt] (0,0) rectangle (14.2,8.3);\node at (7.1,7.700000000000001) {左右点不必同高，但连线的中点始终是中心};
\begin{scope}[shift={(1.491304347826087,2.4166666666666665)},x=1.3043478260869565cm,y=0.5111111111111111cm]\draw[->,gray] (-0.3,0)--(4.3,0) node[right]{$x$};\draw[->,gray] (0,-1.5)--(0,7.5) node[above]{$y$};\draw[blue,thick,solid,domain=0.4:3.6,samples=180] plot(\x,{(\x-2)^3+3});\fill[orange!85!black] (1,2) circle(2.2pt);\fill[orange!85!black] (3,4) circle(2.2pt);\draw[orange!85!black] (1,2)--(3,4);\draw[gray,dashed] (2,-1.2)--(2,7.2);\draw[gray,dashed] (-.2,3)--(4.2,3);\fill[violet] (2,3) circle(2.6pt);\end{scope}\node[anchor=west,align=left,text width=5.0cm] at (8.4,4.0) {$f=(x-2)^3+3$\\[9pt]$P(3,4),\ Q(1,2)$\\[9pt]$\frac{3+1}{2}=2,\quad\frac{4+2}{2}=3$\\[9pt]中点 $C(2,3)$\\[9pt]$f(2+t)+f(2-t)=6$};
\node[text width=12.6cm,align=center] at (7.1,.65) {这条具体曲线示范一般公式；严格判定仍要求所有点都有对应点。};
\end{tikzpicture}\end{center}



取一个数 $t$，从中心横坐标 $a$ 向右走 $t$，得到 $a+t$；向左走同样多，得到 $a-t$。
这两个横坐标天然以 $a$ 为中间数。对应的两点是
\[
P\bigl(a+t,f(a+t)\bigr),\qquad Q\bigl(a-t,f(a-t)\bigr).
\]
若它们的中点正好是 $C(a,b)$，纵坐标的平均数必须等于 $b$：
\[
\frac{f(a+t)+f(a-t)}{2}=b
\quad\Longrightarrow\quad
\boxed{f(a+t)+f(a-t)=2b}.
\]
\textbf{记忆方式：}横坐标在 $a$ 两边“等距离”，纵坐标在 $b$ 两边“等高度、相反方向”。

\medskip

\end{explanationbox}
\clearpage
\begin{explanationbox}
\textbf{\textcolor{CExplanation}{三、为什么还要检查定义域？}}
\begin{center}\begin{tikzpicture}[font=\small,>=stealth,line join=round]
\draw[gray!40,rounded corners=3pt] (0,0) rectangle (14.2,8.3);\node at (7.1,7.700000000000001) {缺少对应点：半条直线不关于原点对称};
\begin{scope}[shift={(4.1,3.9499999999999997)},x=1.0cm,y=0.7666666666666666cm]\draw[->,gray] (-3,0)--(3,0) node[right]{$x$};\draw[->,gray] (0,-3)--(0,3) node[above]{$y$};\draw[blue,thick,solid,domain=0:2.7,samples=180] plot(\x,{\x});\fill[orange!85!black] (0,0) circle(2.2pt);\fill[orange!85!black] (1,1) circle(2.2pt);\end{scope}\node[anchor=west,align=left,text width=5.0cm] at (8.4,4.0) {$f=x,\quad D=[0,+\infty)$\\[9pt]图上有 $(1,1)$\\[9pt]没有 $(-1,-1)$\\[9pt]只验证 $t=0$ 无法判定};
\node[text width=12.6cm,align=center] at (7.1,.65) {原点在图上也不够；必须让所有点转半圈后仍在原图像上。};
\end{tikzpicture}\end{center}



要转半圈后还能落在\textbf{原图像上}，原来的每个点都得有对应的另一点。
如果只有右半边的曲线，左半边根本没画出来，即使个别点碰巧符合等式，也不能叫整幅图像中心对称。
所以检查顺序是：\textbf{先看横坐标成不成对，再看函数值相加是不是同一个常数。}

\medskip

\end{explanationbox}
\clearpage
\begin{explanationbox}
\textbf{\textcolor{CExplanation}{四、把新知识接到已经学过的奇函数}}
\begin{center}\begin{tikzpicture}[font=\small,>=stealth,line join=round]
\draw[gray!40,rounded corners=3pt] (0,0) rectangle (14.2,8.5);\node at (7.1,7.9) {平移证明：把中心搬到原点，就得到奇函数};
\node at (2.9000000000000004,6.9){$f(x)=(x-2)^3+3$};\begin{scope}[shift={(0.7,2.9000000000000004)},x=1.1cm,y=.4cm]\draw[->,gray] (0,0)--(4,0);\draw[->,gray] (0,-2)--(0,8);\draw[blue,thick,domain=0.44999999999999996:3.55,samples=140] plot(\x,{(\x-2)^3+3});\fill[violet] (2,3) circle(2.4pt);\end{scope}\node at (10.100000000000001,6.9){$g(t)=f(2+t)-3=t^3$};\begin{scope}[shift={(10.100000000000001,4.1)},x=1.1cm,y=.4cm]\draw[->,gray] (-2,0)--(2,0);\draw[->,gray] (0,-5)--(0,5);\draw[blue,thick,domain=-1.55:1.55,samples=140] plot(\x,{(\x)^3});\fill[violet] (0,0) circle(2.4pt);\end{scope}\draw[->,orange!85!black,thick] (5.8,4.2)--(7.4,4.2);\node at (7.1,1.4){向左移2、向下移3；两个坐标方向分别使用相同单位长度。};
\node[text width=12.6cm,align=center] at (7.1,.65) {移动后横坐标为 $t=x-2$，高度为 $g(t)=f(x)-3$。};
\end{tikzpicture}\end{center}



奇函数满足 $f(-t)=-f(t)$，意味着 $(t,f(t))$ 与 $(-t,-f(t))$ 围绕\textbf{原点}成对。
如果中心换成 $(a,b)$，就先把图像向左移 $a$、向下移 $b$：
\[
 g(t)=f(a+t)-b,\qquad g(-t)=-g(t).
\]
这不是一套陌生规则，而是\textbf{奇函数的原点对称，经过平移后的同一个道理}。
例如 $y=x^3$ 关于 $(0,0)$ 对称，$y=(x-2)^3+3$ 则关于 $(2,3)$ 对称。


\medskip

\end{explanationbox}
\clearpage
\begin{explanationbox}
\textbf{\textcolor{CExplanation}{五、三个“眼前一亮”的关联}}
\begin{center}\begin{tikzpicture}[font=\small,>=stealth,line join=round]
\draw[gray!40,rounded corners=3pt] (0,0) rectangle (14.2,8.3);\node at (7.1,7.700000000000001) {中心可以不在图上：两支曲线仍绕它成对};
\begin{scope}[shift={(2.1,2.4166666666666665)},x=1.0cm,y=0.5111111111111111cm]\draw[->,gray] (-1,0)--(5,0) node[right]{$x$};\draw[->,gray] (0,-1.5)--(0,7.5) node[above]{$y$};\draw[blue,thick,solid,domain=-0.5:1.7,samples=180] plot(\x,{1/(\x-2)+3});\draw[blue,thick,solid,domain=2.3:4.5,samples=180] plot(\x,{1/(\x-2)+3});\fill[orange!85!black] (1,2) circle(2.2pt);\fill[orange!85!black] (3,4) circle(2.2pt);\draw[gray,dashed] (2,-1.4)--(2,7.4);\draw[gray,dashed] (-.9,3)--(4.9,3);\draw[violet,fill=white] (2,3) circle(3pt);\end{scope}\node[anchor=west,align=left,text width=5.0cm] at (8.4,4.0) {$f=\frac1{x-2}+3$\\[9pt]中心 $(2,3)$\\[9pt]$D=\mathbb R\setminus\{2\}$\\[9pt]$f(1)=2,\quad f(3)=4$\\[9pt]但 $f(2)$ 不存在};
\node[text width=12.6cm,align=center] at (7.1,.65) {紫色空心定位点不属于函数图像；两条虚线是渐近线，不是曲线。};
\end{tikzpicture}\end{center}


\begin{enumerate}[leftmargin=*,itemsep=0.5em]
\item \textbf{中心未必在曲线上：} $f(x)=\dfrac{1}{x-2}+3$ 的图像关于 $(2,3)$ 对称，
但 $x=2$ 处没有函数值；中心像两支曲线之间的“空心定位点”。


\begin{center}\begin{tikzpicture}[font=\small,>=stealth,line join=round]
\draw[gray!40,rounded corners=3pt] (0,0) rectangle (14.2,8.5);\node at (7.1,7.9) {同样取中间横坐标2：轴对称看同高，点对称看平均高度};
\node at (2.4,6.9) {轴：$f=(x-2)^2+3$};\begin{scope}[shift={(0.7,3.0500000000000003)},x=0.85cm,y=0.45cm]\draw[->,gray] (0,0)--(4,0) node[right]{$x$};\draw[->,gray] (0,-1)--(0,7) node[above]{$y$};\draw[blue,thick,solid,domain=0.5:3.5,samples=160] plot(\x,{(\x-2)^2+3});\fill[orange!85!black] (1,4) circle(2pt);\fill[orange!85!black] (3,4) circle(2pt);\end{scope}\node[align=center,text width=4cm] at (2.4,1.7) {$f(1)=f(3)=4$};\node at (7.050000000000001,6.9) {中心：$f=(x-2)^3+3$};\begin{scope}[shift={(5.3500000000000005,3.0500000000000003)},x=0.85cm,y=0.45cm]\draw[->,gray] (0,0)--(4,0) node[right]{$x$};\draw[->,gray] (0,-1)--(0,7) node[above]{$y$};\draw[violet,thick,solid,domain=0.5:3.5,samples=160] plot(\x,{(\x-2)^3+3});\fill[orange!85!black] (1,2) circle(2pt);\fill[orange!85!black] (3,4) circle(2pt);\end{scope}\node[align=center,text width=4cm] at (7.050000000000001,1.7) {$f(1)+f(3)=2+4=6$};\node at (11.7,6.9) {平移后：$g=t^3$};\begin{scope}[shift={(11.7,4.4)},x=0.85cm,y=0.45cm]\draw[->,gray] (-2,0)--(2,0) node[right]{$t$};\draw[->,gray] (0,-4)--(0,4) node[above]{$y$};\draw[orange!85!black,thick,solid,domain=-1.5:1.5,samples=160] plot(\x,{(\x)^3});\fill[orange!85!black] (-1,-1) circle(2pt);\fill[orange!85!black] (1,1) circle(2pt);\end{scope}\node[align=center,text width=4cm] at (11.7,1.7) {$g(-1)=-g(1)$};
\node[text width=12.6cm,align=center] at (7.1,.65) {前两图坐标范围相同；第三图是平移后的坐标范围，纵向显示 $[-4,4]$。};
\end{tikzpicture}\end{center}

\item \textbf{一眼看出奇偶性的区别：}轴对称常见 $f(a+t)=f(a-t)$，两边\textbf{相等}；
点对称常见 $f(a+t)+f(a-t)=2b$，两边\textbf{合起来平均是 $b$}。

\begin{center}\begin{tikzpicture}[font=\small,>=stealth,line join=round]
\draw[gray!40,rounded corners=3pt] (0,0) rectangle (14.2,8.3);\node at (7.1,7.700000000000001) {同一高度的两个中心，半转两次得到水平平移};
\begin{scope}[shift={(2.463636363636364,3.9499999999999997)},x=0.5454545454545454cm,y=1.5333333333333332cm]\draw[->,gray] (-2.5,0)--(8.5,0) node[right]{$x$};\draw[->,gray] (0,-1.5)--(0,1.5) node[above]{$y$};\draw[blue,thick,solid,domain=-2:8,samples=180] plot(\x,{sin(deg(pi*\x/3))});\fill[orange!85!black] (1,0.8660254037844386) circle(2.2pt);\fill[orange!85!black] (-1,-0.8660254037844386) circle(2.2pt);\fill[orange!85!black] (7,0.8660254037844386) circle(2.2pt);\end{scope}\node[anchor=west,align=left,text width=5.0cm] at (8.4,4.0) {图取 $f=\sin\frac{\pi x}{3}$\\[9pt]中心 $(0,0)$、$(3,0)$\\[9pt]$1\to-1\to7$\\[9pt]高度依次 $\frac{\sqrt3}2,-\frac{\sqrt3}2,\frac{\sqrt3}2$\\[9pt]相差 $6$，回到同高};
\node[text width=12.6cm,align=center] at (7.1,.65) {正文对任意函数的配对关系推出 $f(x+6)=f(x)$；6只保证是一个周期。};
\end{tikzpicture}\end{center}

\item \textbf{从两个中心推出周期：}如果同一图像分别关于 $(a,b)$ 和 $(c,b)$ 对称（$a\ne c$），
连续做两次“绕中心转半圈”，结果等于左右平移 $2(c-a)$，因此 $2|c-a|$ 是一个周期。
注意两个中心必须\textbf{一样高}；若纵坐标不同，就不能这样直接说有周期。
\end{enumerate}

\medskip

\end{explanationbox}
\clearpage
\begin{explanationbox}
\textbf{\textcolor{CExplanation}{六、考场先想什么？}}
\begin{center}\begin{tikzpicture}[font=\small,>=stealth,line join=round]
\draw[gray!40,rounded corners=3pt] (0,0) rectangle (14.2,8.3);\node at (7.1,7.700000000000001) {例3：先找平均高度，再补另一侧的函数值};
\begin{scope}[shift={(1,1.8)},x=1cm,y=.65cm]\draw[->,gray] (0,0)--(7,0) node[right]{$x$};\draw[->,gray] (0,0)--(0,8);\draw[blue,thick] (1,7)--(5,3);\foreach \X/\Y in {1/7,5/3}{\fill[orange!85!black] (\X,\Y) circle(2.4pt);}\fill[violet] (3,5) circle(2.6pt);\node[above right=7pt] at (1,7){$(1,7)$};\node[right=7pt] at (5,3){$(5,3)$};\node[below left=7pt,violet] at (3,5){$(3,5)$};\end{scope}\node[anchor=west,align=left] at (9,4.5){两输入和为6\\[12pt]两输出和为10\\[12pt]平均得中心 $(3,5)$};
\node[text width=12.6cm,align=center] at (7.1,.65) {已知 $f(1)=7$，对应 $f(5)=10-7=3$；只画题目保证的一对点。};
\end{tikzpicture}\end{center}


给出形如 $f(x)+f(6-x)=10$：不要先求解析式，先看两个输入的中间数是 $3$，
两个输出的中间数是 $5$，对称中心候选是 $(3,5)$；\textbf{最后别忘了定义域}。
\end{explanationbox}
